% !PRINTED 235
\weberpage{235}
%<<CONT
und nach (4) kann diese Gleichung auch so dargestellt werden:
\weq{7}{m\,(a\,u + b\,v) = (a'\,u + b'\,v)^2 - c\,(a\,u + b\,v)^2.}
Man setze nun
\weq{8}{y = \frac{a'\,u + b'\,v}{a\,u + b\,v},}
und bilde die Gleichung f\"unften Grades, deren Wurzeln die f\"unf Werthe von $y$ sind:
\weq{9}{y^5 + c_1\,y^4 + c_2\,y^3 + c_3\,y^2 + c_4\,y + c_5 = 0;}
darin lassen sich die Co\"efficienten $c$ auf folgende Weise n\"aher bestimmen.

Aus (8) leiten wir ab
\[
y + \sqrt{c} = \frac{a'\,u + b'\,v + \sqrt{c}\,(a\,u + b\,v)}{a\,u + b\,v}
\]
und daraus nach (7)
\[
y + \sqrt{c} = \frac{m}{a'\,u + b'\,v - \sqrt{c}\,(a\,u + b\,v)},
\]
\[
\frac{m}{y + \sqrt{c}} = a'\,u + b'\,v - \sqrt{c}\,(a\,u + b\,v).
\]
Daraus folgt aber nach (3)
\[
S\left[\frac{1}{y + \sqrt{c}}\right] = 0, \quad S\left[\frac{1}{(y + \sqrt{c})^2}\right] = 0.
\]
Wenn wir also aus (9) die Gleichung ableiten, deren Wurzeln die Werthe
\[
\xi = \frac{1}{y + \sqrt{c}}
\]
sind, also
\[
y = \frac{1 - \xi\,\sqrt{c}}{\xi}
\]
setzen, so muss eine Gleichung f\"ur $\xi$ entstehen, in der die Co\"efficienten von $\xi^4$ und $\xi^3$ verschwinden, und zwar welches Zeichen wir auch der Quadratwurzel $\sqrt{c}$ geben. Dadurch erh\"alt man vier Gleichungen zwischen den Co\"efficienten $c_\nu$ der Gleichung (9).

Die Gleichung f\"ur $\xi$ wird n\"amlich
\[
(1 - \xi\,\sqrt{c})^5 + c_1\,\xi\,(1 - \xi\,\sqrt{c})^4 + c_2\,\xi^2\,(1 - \xi\,\sqrt{c})^3 + c_3\,\xi^3\,(1 - \xi\,\sqrt{c})^2 + c_4\,\xi^4\,(1 - \xi\,\sqrt{c}) + c_5\,\xi^5 = 0.
\]
Setzt man hierin die Co\"efficienten von $\xi^4$ und $\xi^3$ gleich 0, so folgt
%CONT>>
